%% text-sharing.tex -- shared closing section
\subsubsection{Building and sharing}

The package is \texttt{florilegium.sty} with nine Lua files beside it: \texttt{ink.lua}
(geometry, color and PDF output), \texttt{flora.lua} (the primitives),
\texttt{authored.lua} and \texttt{grow.lua} (the two compositions), \texttt{initial.lua}
(the decorated initials), two themes and the figures one of them draws from
(\texttt{theme-geometric.lua}: the same plant cut from flat shapes;
\texttt{theme-math.lua}: figures, symbols and equations cascading down the margin;
\texttt{figures.lua}) and \texttt{florilegium.lua} (the glue to \TeX). It needs Lua\LaTeX\ and one compilation, at
about a second per page. To share it, the three routes that work are, in increasing order
of effort: paste the files into a blank Overleaf project with the compiler set to
Lua\LaTeX\ and share the link; put them in a Git repository with a continuous integration
job that compiles the PDFs on every push and attaches them to tagged releases; or publish
the package so that any document can \texttt{\textbackslash usepackage} the garden and
choose its own seed, column and hue.

\begin{remark}
Two mistakes cost time and are worth recording. A literal written into the page bypasses
TikZ, so it also bypasses TikZ's transformations: a primitive called inside
\verb|\begin{scope}[shift=...,rotate=...]| must read pgf's current matrix itself (it is in
\verb|\pgf@pt@aa| to \verb|\pgf@pt@bb| and the dimensions \verb|\pgf@pt@x|,
\verb|\pgf@pt@y|) and apply it in Lua. And a Lua function that returns two values, placed
in the middle of a table constructor such as \verb|{x0, y0, f(t), x1, y1}|, contributes only
its first value; the rest of the table shifts, and the error surfaces three functions
later as arithmetic on \texttt{nil}.
\end{remark}
